\begin{tikzpicture}[font=\small, >=latex]
  \tikzset{bus/.style={circle, draw=blue!65!black, fill=blue!10, minimum size=13mm, align=center},
    line/.style={thick, draw=gray!80!black}}
  \node[bus] (a) at (0,0) {1\\$\theta_1$};
  \node[bus] (b) at (3.3,1.35) {2\\$\theta_2$};
  \node[bus] (c) at (3.3,-1.35) {3\\$\theta_3$};
  \draw[line] (a) -- node[above, sloped, font=\scriptsize] {$f_{12}=b_{12}(\theta_1-\theta_2)$} (b);
  \draw[line] (a) -- node[below, sloped, font=\scriptsize] {$f_{13}=b_{13}(\theta_1-\theta_3)$} (c);
  \draw[line] (b) -- node[right, font=\scriptsize] {$f_{23}$} (c);
  \node[draw, rounded corners=2pt, fill=orange!12, draw=orange!70!black, align=left, text width=47mm] (eq) at (7.25,0) {DCOPF teaching fixture\\[1mm]
    minimize generation cost\\
    subject to nodal balance\\
    $|f_{ij}|\leq F_{ij}$\\
    fix one reference angle};
  \draw[->, thick] (b) -- (eq);
  \draw[->, thick] (c) -- (eq);
  \node[below=9mm of a, align=center, text width=105mm] {A three-bus fixture makes the DCOPF equations inspectable before using the ACTIVSg2000 planning model. The student CSV is an input to the larger model, not an hourly dispatch table. Student preview can check residuals; the evaluator remains authoritative.};
\end{tikzpicture}
